\documentclass{article}
\usepackage{amsmath}

\title{Worked Example: The Case $p=25$}
\author{Matteo Calvigioni}
\date{}

\begin{document}

\maketitle

\section*{Introduction}

We illustrate the congruence
\[
2S(n)\equiv n \pmod{p}
\]
in the concrete case $p=25$, using the number
\[
n=57342.
\]

\section*{Digits and digit sum}

We write
\[
n=57342 = \sum_{i=0}^4 d_i 10^i,
\]
with
\[
d_0=2,\quad d_1=4,\quad d_2=3,\quad d_3=7,\quad d_4=5.
\]

The digit sum is
\[
S(n)=5+7+3+4+2=21,
\]
hence
\[
2S(n)=42.
\]

\section*{Reduction modulo $25$}

Since $25 \mid 100$, we have
\[
10^2 \equiv 0 \pmod{25},
\]
and therefore
\[
n \equiv d_0 + 10d_1 \pmod{25}.
\]

Thus
\[
57342 \equiv 2 + 10\cdot 4 = 42 \equiv 17 \pmod{25}.
\]

At the same time,
\[
2S(n)=42 \equiv 17 \pmod{25}.
\]

Therefore
\[
2S(57342)\equiv 57342 \pmod{25}.
\]

\section*{Verification via the weighted digit formula}

We also consider
\[
\sum_{i=0}^m d_i(2-10^i)\equiv 0 \pmod{25}.
\]

Since
\[
10^0\equiv 1,\quad 10^1\equiv 10,\quad 10^i\equiv 0 \pmod{25} \ \text{for } i\ge 2,
\]
the congruence becomes
\[
d_0 + 17d_1 + 2\sum_{i=2}^m d_i \equiv 0 \pmod{25}.
\]

For $n=57342$ we have
\[
d_0=2,\quad d_1=4,\quad d_2+d_3+d_4=3+7+5=15,
\]
so
\[
2 + 17\cdot 4 + 2\cdot 15 = 2 + 68 + 30 = 100 \equiv 0 \pmod{25}.
\]

\section*{Conclusion}

The number $n=57342$ satisfies
\[
2S(n)\equiv n \pmod{25}.
\]

This example illustrates the original motivating identity
\[
2S(n) = p + (n \bmod p)
\]
in the case $p=25$.

\end{document}
